% =============================================================================
% Chapter 13: Security and Privacy — ML Systems Lecture Slides (Vol II)
% =============================================================================
% IMAGES:
%   - security-vs-privacy.pdf      - attack-surface-layers.pdf
%   - supply-chain-attack.pdf      - ml-lifecycle-threats.pdf
%   - model-theft-economics.pdf    - data-poisoning.pdf
%   - differential-privacy.pdf     - defense-stack.pdf
%   - hw-security-stack.pdf        - implementation-roadmap.pdf
% =============================================================================
\documentclass[aspectratio=169, 12pt]{beamer}
\usepackage{../../assets/beamerthememlsys}

\mlsyssetup{
  volume       = {Volume II},
  chapter      = {Chapter 13},
  logo         = {../../assets/img/logo-mlsysbook.png},
  instlogo     = {../../assets/img/logo-harvard.png},
  chaptertitle = {Security and Privacy},
}

% --- Fonts ---
\usepackage{fontspec}
\setsansfont{Helvetica Neue}[
  BoldFont={Helvetica Neue Bold},
  ItalicFont={Helvetica Neue Italic},
  BoldItalicFont={Helvetica Neue Bold Italic},
]
\setmonofont{JetBrains Mono}[Scale=0.85]

% --- Packages ---
\usepackage{booktabs}
\usepackage{amsmath}

% --- Image paths ---
\graphicspath{
  {images/}
}

% --- Chapter-specific macros ---
\newcommand{\DP}{Differential Privacy}

% --- Helper: safe image include ---
\newcommand{\safeimg}[2][width=\textwidth,keepaspectratio]{%
  \IfFileExists{images/#2}{\includegraphics[#1]{#2}}{%
    \IfFileExists{#2}{\includegraphics[#1]{#2}}{%
      \fbox{\parbox[c][2.5cm][c]{0.85\linewidth}{%
        \centering\footnotesize\textcolor{midgray}{[Missing image]}}}%
    }%
  }%
}

% --- Section count for navigation ---
\setcounter{mlsystotalsections}{8}

\title{Security and Privacy}
\author{Vijay Janapa Reddi}
\institute{Harvard University}
\date{}

\begin{document}

% =============================================================================
% TITLE SLIDE
% =============================================================================
\mlsystitle{Security \& Privacy}{Defending the ML Fleet}{cover_security_privacy.png}

% =============================================================================
% LEARNING OBJECTIVES
% =============================================================================
\begin{frame}{Learning Objectives}
\note{[2 min] Walk through objectives. Emphasize that this chapter is about threats
unique to ML --- not just traditional IT security. Ask: ``What makes securing an ML
system fundamentally harder than securing a web application?''}

\small
\begin{enumerate}
  \item Distinguish \textbf{security} from \textbf{privacy} in ML through formal definitions and \textbf{threat models}
  \item Extract attack patterns from historical breaches (\textbf{Stuxnet}, \textbf{Jeep Cherokee}, \textbf{Mirai})
  \item Analyze ML-specific attacks: \textbf{model theft}, \textbf{data poisoning}, \textbf{adversarial examples}
  \item Implement \textbf{differential privacy} with privacy budgets and accuracy trade-offs
  \item Design \textbf{layered defense} architectures from data to hardware
  \item Evaluate \textbf{hardware security primitives} (TEEs, HSMs, PUFs) with overhead analysis
  \item Apply the \textbf{three-phase roadmap} for context-appropriate security
\end{enumerate}

\end{frame}

\begin{frame}{Visual Language}
\note{[1 min] Explain the semantic color system used throughout the course.
These colors are consistent across all diagrams and slides.}

\small
Throughout this course, colors carry meaning:

\vspace{0.3cm}
\begin{columns}[T]
  \begin{column}{0.45\textwidth}
    \begin{mlsyscard}{computestroke}
    \textbf{Blue} --- Compute / Processing\\
    {\footnotesize GPU ops, forward/backward pass, inference}
    \end{mlsyscard}
    \vspace{0.15cm}
    \begin{mlsyscard}{datastroke}
    \textbf{Green} --- Data / Memory\\
    {\footnotesize Data flow, caches, healthy paths}
    \end{mlsyscard}
  \end{column}
  \begin{column}{0.45\textwidth}
    \begin{mlsyscard}{routingstroke}
    \textbf{Orange} --- Routing / Scheduling\\
    {\footnotesize Load balancers, batch windows}
    \end{mlsyscard}
    \vspace{0.15cm}
    \begin{mlsyscard}{errorstroke}
    \textbf{Red} --- Error / Cost / Bottleneck\\
    {\footnotesize Loss, decode phase, waste}
    \end{mlsyscard}
  \end{column}
\end{columns}
\end{frame}



% =============================================================================
\section{Attack Surface}
% =============================================================================

\begin{frame}{The Expanded ML Attack Surface}
\note{[3 min] Key insight: ML systems do not merely store data --- they learn from it
and embed patterns into persistent weights. A database breach steals records; an ML
breach can reverse-engineer the training data. Ask: ``When a model memorizes a patient
record, who is responsible for that privacy violation?''
Common error: students think traditional IT security is sufficient.}

\footnotesize
\begin{columns}[T]
  \begin{column}{0.55\textwidth}
    ML systems create \alert{fundamentally new} attack surfaces:

    \vspace{0.1cm}
    \begin{itemize}\setlength\itemsep{1pt}
      \item Models \textbf{memorize} training data in weights
      \item Gradient sync creates \textbf{leakage channels}
      \item APIs expose \textbf{side channels} (logits, confidence)
      \item Multi-tenant serving enables \textbf{cross-tenant attacks}
    \end{itemize}

    \vspace{0.1cm}
    \mlsysalert{Key:}{Traditional firewalls cannot defend against adversarial inputs --- these exploit the model's \textbf{learned functions}.}
  \end{column}
  \begin{column}{0.42\textwidth}
    \safeimg[width=\textwidth,height=4.5cm,keepaspectratio]{attack-surface-layers.pdf}
  \end{column}
\end{columns}

\end{frame}

\begin{frame}{Security vs.\ Privacy: Distinct Concerns}
\note{[3 min] Walk through the comparison. Security = adversarial manipulation;
Privacy = unauthorized inference from learned representations. Key insight: they
can be in tension --- robust models are MORE susceptible to membership inference.
Ask: ``Can you have perfect security without privacy? Privacy without security?''}

% --- Layout: FULL-WIDTH diagram ---
\centering
\safeimg[width=0.92\textwidth,height=4.5cm,keepaspectratio]{security-vs-privacy.pdf}

\vspace{0.1cm}
\small
\mlsysalert{Key distinction:}{Security requires \emph{worst-case} analysis; Privacy requires \emph{statistical guarantees}.}

\end{frame}

% --- ACTIVE LEARNING 1: Predict ---
\begin{frame}{Predict: What Makes ML Security Different?}
\note{[2 min] Prediction exercise. Give students 60 seconds. Do NOT reveal the
answer yet. The key insight is that ML adds an algorithmic attack surface
orthogonal to infrastructure. Ask 2--3 students to share.}

\centering
\vspace{1.0cm}
{\Large\bfseries Think--Write--Share}

\vspace{0.6cm}
{\large A properly authenticated API request containing an\\
adversarial input bypasses all network-layer defenses.\\[0.3cm]
\alert{Why can't traditional IT security protect ML systems?}}

\vspace{0.5cm}
{\normalsize Write down 2 reasons. \textcolor{midgray}{(60 seconds)}}

\vspace{0.4cm}
{\small\textcolor{lightgray}{Compare your answer with a neighbor.}}

\end{frame}

% =============================================================================
\section{Historical Lessons}
% =============================================================================

\begin{frame}{Learning from Breaches: Stuxnet}
\note{[3 min] Stuxnet (2010): chained 4 zero-day exploits to reach air-gapped PLCs.
Key insight: it altered centrifuge speeds while reporting normal telemetry ---
silent manipulation beats loud destruction. The ML parallel: a poisoned model
shifts decision boundaries silently. Ask: ``What is the ML equivalent of fake
telemetry?'' Answer: a model with 95\%+ clean accuracy but a hidden backdoor.}

% --- Layout: FULL-WIDTH diagram ---
\centering
\safeimg[width=0.92\textwidth,height=4.5cm,keepaspectratio]{supply-chain-attack.pdf}

\vspace{0.1cm}
\small
\mlsysinsight{Stuxnet lesson:}{Silent manipulation of learned parameters is more devastating than crashing the system.}

\end{frame}

\begin{frame}{Jeep Cherokee \& Mirai: More Parallels}
\note{[2 min] Two more case studies, cover quickly. Jeep Cherokee (2015):
insufficient CAN bus isolation let attackers control steering from the infotainment
system. ML parallel: multi-tenant GPU sharing without hardware isolation. Mirai
(2016): 600K devices with default credentials weaponized into a botnet. ML parallel:
edge ML fleet with unpatched firmware. If short on time, focus on the table.}

\footnotesize
\renewcommand{\arraystretch}{1.15}
\begin{tabular}{@{}lp{3.5cm}p{3.5cm}p{3.5cm}@{}}
  \toprule
  \textbf{Incident} & \textbf{Vulnerability} & \textbf{ML Parallel} & \textbf{Quantitative Impact} \\
  \midrule
  \textbf{Stuxnet} (2010) & Supply chain compromise via 4 zero-day exploits & Poisoned PyPI packages, backdoored model weights & Multi-million \$ weapon \\
  \textbf{Jeep Cherokee} (2015) & No CAN bus isolation between infotainment and control & Multi-tenant GPU without MIG isolation & 1.4M vehicle recall \\
  \textbf{Mirai} (2016) & Default credentials on 600K IoT devices & Unpatched edge ML firmware across fleet & 1 Tbps DDoS \\
  \bottomrule
\end{tabular}

\vspace{0.15cm}
\mlsysalert{Pattern:}{Every attack vector from these incidents has a direct analog in modern ML deployments.}

\end{frame}

% =============================================================================
\section{Attack Vectors}
% =============================================================================

\begin{frame}{Threats Across the ML Lifecycle}
\note{[3 min] Walk through lifecycle stages. Each stage has distinct threats and
requires distinct defenses. Key insight: different threats need different defenses
at different layers. Ask: ``Which lifecycle stage is most vulnerable?''
Common error: students focus only on inference-time attacks.}

% --- Layout: FULL-WIDTH diagram ---
\centering
\safeimg[width=0.88\textwidth,height=4.8cm,keepaspectratio]{ml-lifecycle-threats.pdf}

\vspace{0.1cm}
\small
\mlsysconcept{Lifecycle mapping:}{Match defenses to the stage where the threat operates.}

\end{frame}

\begin{frame}{Model Theft: 100{,}000$\times$ Cost Asymmetry}
\note{[3 min] Walk through the economics. A GPT-4-class model costs \$100M to
train but can be extracted for \$1,000 via black-box queries. 10K--100K queries
achieve 90\% functional accuracy. API rate limits of 100K/day make this trivially
achievable. Defenses: rate limiting, output perturbation, top-k filtering.
Ask: ``If you were the attacker, how would you minimize detection?''}

% --- Layout: FULL-WIDTH diagram ---
\centering
\safeimg[width=0.88\textwidth,height=4.8cm,keepaspectratio]{model-theft-economics.pdf}

\vspace{0.1cm}
\small
\mlsysalert{Every API response leaks information:}{Confidence scores and logits form a side channel for extraction.}

\end{frame}

% --- ACTIVE LEARNING: Micro-Retrieval Cue ---
\begin{frame}{Quick Check}
\note{[1 min] Answer: 100,000x --- 100M dollar training vs 1K dollar extraction.}

\centering
\vspace{1.0cm}
{\Large\bfseries Quick Check}

\vspace{0.6cm}
{\large What is the cost asymmetry in model theft?}

\vspace{0.5cm}
{\normalsize\textcolor{midgray}{15 seconds --- then I will cold-call.}}

\end{frame}



\begin{frame}{Data Poisoning: Extreme Leverage}
\note{[3 min] The most asymmetric attack. 0.1\% corruption (1K out of 1M samples)
degrades accuracy 10--50\%. Backdoor attacks are even more efficient: 0.01\%
creates models with 95\%+ clean accuracy but 90\%+ attack success on triggers.
Backdoors survive retraining and transfer learning.
Ask: ``How would you detect 1,000 poisoned samples in 1 million?''
Common error: students think only large-scale corruption matters.}

\small
\begin{columns}[T]
  \begin{column}{0.55\textwidth}
    \safeimg[width=\textwidth,height=5.5cm,keepaspectratio]{data-poisoning.pdf}
  \end{column}
  \begin{column}{0.42\textwidth}
    \textbf{Attack economics favor adversaries:}

    \vspace{0.15cm}
    \textcolor{errorstroke}{\textbf{Poisoning:}} 0.1\% of data\\
    {\footnotesize $\to$ 10--50\% accuracy drop}

    \vspace{0.1cm}
    \textcolor{errorstroke}{\textbf{Backdoors:}} 0.01\% of data\\
    {\footnotesize $\to$ 90\%+ attack success rate}\\
    {\footnotesize $\to$ 95\%+ clean accuracy preserved}

    \vspace{0.15cm}
    \begin{mlsyscard}{errorstroke}
    {\footnotesize Backdoors \textbf{persist} through retraining and transfer learning. ``Clean enough'' data does not mean ``safe enough'' models.}
    \end{mlsyscard}
  \end{column}
\end{columns}

\end{frame}

\begin{frame}{Adversarial Examples \& Prompt Injection}
\note{[2 min] Two runtime attacks. Adversarial examples: imperceptible perturbations
(less than 0.01\% pixel change) cause misclassification. Transfer across architectures
with 60--80\% success. Prompt injection: exploits LLMs' inability to distinguish
instruction from content. Direct injection overrides system prompts; indirect
injection manipulates retrieved context. Both bypass all network-layer defenses.
If short: focus on prompt injection as the newer, more relevant threat.}

\small
\begin{columns}[T]
  \begin{column}{0.48\textwidth}
    \begin{mlsyscard}{errorstroke}
    \textbf{Adversarial Examples}\\[0.1cm]
    {\footnotesize
    $<$0.01\% pixel change $\to$ misclassification\\[0.05cm]
    Transfer across architectures: 60--80\%\\[0.05cm]
    Stop sign $+$ tape $\to$ ``speed limit'' \\[0.05cm]
    Exploits shared geometric properties of decision boundaries}
    \end{mlsyscard}
  \end{column}
  \begin{column}{0.48\textwidth}
    \begin{mlsyscard}{crimson}
    \textbf{Prompt Injection (LLM-specific)}\\[0.1cm]
    {\footnotesize
    \textbf{Direct:} ``Ignore previous instructions\ldots''\\[0.05cm]
    \textbf{Indirect:} Poison retrieved context in RAG\\[0.05cm]
    Exploits inability to distinguish instruction from content\\[0.05cm]
    No buffer overflow --- the model IS the attack surface}
    \end{mlsyscard}
  \end{column}
\end{columns}

\vspace{0.15cm}
\mlsysalert{Both attacks:}{bypass all network-layer defenses --- they exploit the \textbf{learned function} itself.}

\end{frame}

% --- ACTIVE LEARNING 2: Exercise ---
\begin{frame}{Your Turn: Assess the Threat}
\note{[3 min] Give students 90 seconds. Answer: This is a model extraction attack.
Rate limiting alone is insufficient because 100K queries/day = extractable in 1 day.
Need: output perturbation + top-k filtering + query pattern anomaly detection.
The key insight is that rate limits are necessary but not sufficient.}

\small
\begin{columns}[T]
  \begin{column}{0.58\textwidth}
    {\normalsize\bfseries Diagnose the Attack}

    \vspace{0.2cm}
    An API serves a medical diagnosis model. You observe:
    \begin{itemize}
      \item 50K queries/day from one account (limit: 100K)
      \item Queries are systematically varied inputs
      \item Each response includes full confidence scores
    \end{itemize}

    \vspace{0.15cm}
    \textbf{Identify} the attack type and propose 3 defenses.

    {\footnotesize\textcolor{midgray}{(90 seconds --- then compare with a neighbor)}}
  \end{column}
  \begin{column}{0.38\textwidth}
    \pause
    \begin{mlsyscard}{errorstroke}
    \textbf{Solution:}\\[0.1cm]
    {\footnotesize
    \textbf{Attack:} Model extraction\\[0.1cm]
    \textbf{Defenses:}\\
    1. Output perturbation (add noise to logits)\\
    2. Top-k filtering (return only top-3)\\
    3. Query pattern anomaly detection\\[0.1cm]
    Rate limits alone $\neq$ sufficient!
    }
    \end{mlsyscard}
  \end{column}
\end{columns}

\end{frame}

% =============================================================================
\section{Hardware Security}
% =============================================================================

\begin{frame}{Hardware-Level Vulnerabilities}
\note{[3 min] Hardware is the foundation of the trust chain. Six categories:
bugs (Spectre/Meltdown), physical attacks (probing), fault injection (laser/voltage
glitch), side-channels (power/timing/EM), leaky interfaces (DMA, JTAG), counterfeit
chips. Key insight: software security is only as strong as the hardware.
If short: focus on side-channels and TEEs as the two most ML-relevant topics.}

\scriptsize
\renewcommand{\arraystretch}{1.0}
\begin{tabular}{@{}lp{3cm}p{3.5cm}p{2.8cm}@{}}
  \toprule
  \textbf{Category} & \textbf{Mechanism} & \textbf{ML Impact} & \textbf{Defense} \\
  \midrule
  \textbf{HW Bugs} & Spectre/Meltdown & Leak weights from shared cache & Microcode, isolation \\
  \textbf{Physical} & Chip probing & Extract weights from edge devices & Tamper-resistant pkg \\
  \textbf{Fault Inj.} & Laser/voltage glitch & Bit-flip params for backdoor & Error detection \\
  \textbf{Side-Chan.} & Power/timing/EM & Recover keys, infer data & Constant-time code \\
  \textbf{Leaky I/F} & DMA, JTAG ports & Direct access to model weights & Disable in production \\
  \textbf{Counterfeit} & Fake supply chain chips & Reduced security margins & PUFs, audit \\
  \bottomrule
\end{tabular}

\vspace{0.05cm}
\footnotesize
\mlsysalert{Key:}{Software protections are only as secure as the silicon they run on.}

\end{frame}

\begin{frame}{Hardware Security Primitives}
\note{[3 min] Four primitives that anchor ML security. TEEs: encrypted memory
enclaves (5--30\% throughput cost). Secure Boot: chain of trust from silicon to
ML runtime. HSMs: dedicated crypto processors (2--10x latency). PUFs: unclonable
chip fingerprints for edge device authentication. Ask: ``Which primitive would you
deploy first for a multi-tenant ML cloud?'' Answer: TEEs.}

% --- Layout: FULL-WIDTH diagram ---
\centering
\safeimg[width=0.88\textwidth,height=4.2cm,keepaspectratio]{hw-security-stack.pdf}

\vspace{0.1cm}
\footnotesize
\mlsysconcept{Multi-tenant clouds $\to$ TEEs;}{edge fleets $\to$ PUFs; model registries $\to$ HSMs.}

\end{frame}

\begin{frame}{The Tax of Trusted Compute}
\note{[2 min] Worked example: MIG partitioning on H100. Dedicated GPU: 1,000
tokens/sec. After MIG with secure context switching: 850 tokens/sec. That is a
15\% throughput loss --- the ``isolation tax.'' For multi-tenant fleets, this is
the price of preventing one tenant's prompt from leaking into another's response.
Ask: ``When is 15\% throughput loss worth paying?''}

\small
\begin{columns}[T]
  \begin{column}{0.55\textwidth}
    \textbf{Multi-Instance GPU (MIG) Isolation}

    \vspace{0.15cm}
    {\footnotesize
    \renewcommand{\arraystretch}{1.15}
    \begin{tabular}{@{}lr@{}}
      \toprule
      \textbf{Metric} & \textbf{Value} \\
      \midrule
      Dedicated GPU throughput & 1,000 tok/s \\
      After MIG partitioning   & 850 tok/s \\
      \textbf{Isolation Tax}   & \textcolor{errorstroke}{\textbf{15\%}} \\
      \bottomrule
    \end{tabular}
    }

    \vspace{0.15cm}
    \begin{mlsyscard}{routingstroke}
    {\footnotesize Security is a \textbf{capacity drain}. Providing a secure enclave costs 15\% of raw throughput. Multi-tenancy is an economic necessity, but isolation is not free.}
    \end{mlsyscard}
  \end{column}
  \begin{column}{0.42\textwidth}
    \vspace{0.2cm}
    {\footnotesize\textbf{TEE Overhead Ranges:}}

    \vspace{0.1cm}
    {\footnotesize
    \renewcommand{\arraystretch}{1.15}
    \begin{tabular}{@{}lr@{}}
      \toprule
      \textbf{Primitive} & \textbf{Cost} \\
      \midrule
      TEE (SGX/SEV) & 5--30\% \\
      Secure Boot   & Boot time \\
      HSM key ops   & 2--10$\times$ \\
      PUF auth      & Minimal \\
      \bottomrule
    \end{tabular}
    }
  \end{column}
\end{columns}

\end{frame}

% =============================================================================
\section{Differential Privacy}
% =============================================================================

\mlsysfocus{Differential Privacy}{%
$\Pr[\mathcal{M}(D) \in S] \leq e^{\epsilon} \cdot \Pr[\mathcal{M}(D') \in S] + \delta$%
}

\begin{frame}{Differential Privacy: Calibrated Noise}
\note{[3 min] DP adds calibrated noise so no individual's data can be reverse-
engineered from the output. Key parameters: epsilon (privacy budget), sensitivity
(max one record changes output), noise scale = S/epsilon. The trade-off is
quantitative: epsilon=0.1 gives strong privacy but 10--15\% accuracy loss.
Walk through the diagram. Ask: ``Why does DP only work at scale?''
Answer: error per person ~ Noise/N, so larger N dampens the noise.}

% --- Layout: FULL-WIDTH diagram ---
\centering
\safeimg[width=0.92\textwidth,height=4.5cm,keepaspectratio]{differential-privacy.pdf}

\vspace{0.1cm}
\small
\mlsysinsight{Scale matters:}{Error per person $\approx$ Noise$/N$. DP kills utility for small $N$ but works at fleet scale.}

\end{frame}

\begin{frame}{DP Worked Example: Salary Privacy}
\note{[2 min] Concrete numbers. N=1000 employees, salary range \$0--\$200K,
epsilon=1. Sensitivity = \$200K. Noise scale = 200K/1 = \$200K. Error per person
= \$200K/1000 = \$200. For N=100, error = \$2,000 --- DP is impractical.
This is why Apple uses DP for emoji usage (billions of events) but not for
small clinical trials.}

\small
\textbf{Computing average salary with $\epsilon=1.0$ privacy}

\vspace{0.15cm}
\renewcommand{\arraystretch}{1.2}
{\footnotesize
\begin{tabular}{@{}llll@{}}
  \toprule
  \textbf{Step} & \textbf{Parameter} & \textbf{Value} & \textbf{Meaning} \\
  \midrule
  Sensitivity ($S$)    & Max one person changes sum & \$200,000 & Range of salaries \\
  Privacy budget       & $\epsilon$ & 1.0 & Moderate privacy \\
  Noise scale ($b$)    & $S / \epsilon$ & \$200,000 & Laplace noise magnitude \\
  Error ($N$=1,000)    & $b / N$ & \textcolor{datastroke}{\textbf{\$200}} & Acceptable \\
  Error ($N$=100)      & $b / N$ & \textcolor{errorstroke}{\textbf{\$2,000}} & Impractical \\
  \bottomrule
\end{tabular}
}

\vspace{0.15cm}
\begin{mlsyscard}{routingstroke}
{\footnotesize \textbf{The Privacy-Utility Trade-off}: $\epsilon=0.1$ gives strong privacy but 10--15\% accuracy loss. $\epsilon=1.0$ is moderate (5--10\% loss). $\epsilon=10$ is weak but minimal loss.}
\end{mlsyscard}

\end{frame}

\begin{frame}{Privacy Budget Composition}
\note{[2 min] Critical implementation detail: privacy budgets COMPOUND. Training 10
models at epsilon=1 each = total epsilon=10. Monthly retraining for 2 years accumulates
epsilon=24 even if each run targets epsilon=1. Organizations failing to track cumulative
loss exceed guarantees by orders of magnitude. DP-SGD applies to training: clip gradients,
add noise. Cost: 5--10\% accuracy at epsilon=1.
If short: just emphasize the compounding.}

\small
\begin{columns}[T]
  \begin{column}{0.55\textwidth}
    \textbf{Privacy budgets compound:}

    \vspace{0.15cm}
    \begin{itemize}\setlength\itemsep{2pt}
      \item 10 models $\times$ $\epsilon=1.0$ each $=$ total $\epsilon=10$
      \item Monthly retrain $\times$ 24 months $=$ $\epsilon=24$
      \item Hyperparameter tuning adds more budget
    \end{itemize}

    \vspace{0.1cm}
    \begin{mlsyscard}{errorstroke}
    {\footnotesize Organizations failing to track cumulative privacy loss exceed guarantees by \textbf{orders of magnitude} --- violating regulations despite using DP libraries.}
    \end{mlsyscard}
  \end{column}
  \begin{column}{0.42\textwidth}
    {\footnotesize\textbf{DP-SGD for Training:}}
    \vspace{0.1cm}
    \begin{enumerate}\setlength\itemsep{1pt}
      \item[\footnotesize 1.] Clip per-sample gradients
      \item[\footnotesize 2.] Add calibrated Gaussian noise
      \item[\footnotesize 3.] Track privacy budget per step
      \item[\footnotesize 4.] Stop when budget exhausted
    \end{enumerate}

    \vspace{0.15cm}
    {\footnotesize\textbf{Cost:} 5--10\% accuracy loss at $\epsilon \approx 1.0$. Requires 2--5$\times$ more training steps.}
  \end{column}
\end{columns}

\end{frame}

% =============================================================================
\section{Defense Architecture}
% =============================================================================

\begin{frame}{Layered Defense Architecture}
\note{[3 min] Five layers from silicon to governance. Key insight: no single layer
suffices; composition creates graceful degradation. Walk through each layer briefly.
Ask: ``What happens if you skip the hardware root of trust?''
Answer: all software-layer protections can be undermined.}

% --- Layout: FULL-WIDTH diagram ---
\centering
\safeimg[width=0.92\textwidth,height=4.8cm,keepaspectratio]{defense-stack.pdf}

\vspace{0.1cm}
\small
\mlsysconcept{Defense in depth:}{Each layer catches threats the layers above cannot.}

\end{frame}

\begin{frame}{Secure Aggregation for Federated Learning}
\note{[2 min] FL privacy challenge: gradient updates leak data. Gradient inversion
can reconstruct training images from gradient vectors. Membership inference achieves
70--90\% accuracy on FL models. Defense: secure aggregation prevents server from
seeing individual contributions; add DP noise to updates (epsilon approx 6).
The Google Gboard example: next-word prediction trained on millions of phones
without centralizing any user text.}

\footnotesize
\begin{columns}[T]
  \begin{column}{0.55\textwidth}
    \textbf{Federated learning is NOT inherently private:}

    \vspace{0.1cm}
    \begin{itemize}\setlength\itemsep{1pt}
      \item Gradient inversion reconstructs training data
      \item Membership inference: 70--90\% accuracy
      \item Raw gradients $\to$ the server sees everything
    \end{itemize}

    \vspace{0.1cm}
    \mlsysinsight{Required:}{Secure aggregation + DP ($\epsilon \approx 6$) + cryptographic protection.}
  \end{column}
  \begin{column}{0.42\textwidth}
    \vspace{0.2cm}
    \begin{mlsyscard}{crimson}
    \textbf{Google Gboard}\\[0.1cm]
    {\footnotesize Next-word prediction trained on millions of phones. No text leaves device. Secure aggregation + DP noise protect typing patterns.}
    \end{mlsyscard}
  \end{column}
\end{columns}

\end{frame}

% --- ACTIVE LEARNING 3: Discussion ---
\begin{frame}{Discussion: Where Is the Weakest Link?}
\note{[3 min] Turn-and-talk. Students discuss in pairs for 90 seconds. Cold-call
2--3 pairs. The point: a single weak link compromises an otherwise secure system.
Strong API defenses are useless if model weights leak through CI/CD pipelines
or monitoring logs. The answer should be ``it depends on the deployment'' ---
the exercise surfaces that security is a systems property, not a component property.}

\centering
\vspace{0.6cm}
{\large\bfseries Turn and Talk \textcolor{midgray}{(90 seconds)}}

\vspace{0.4cm}
{\large A production ML system has strong API rate limiting,\\
encrypted model storage, and DP in training.\\[0.3cm]
\alert{An attacker still extracts the model. How?}}

\vspace{0.4cm}
\begin{columns}[c]
  \begin{column}{0.19\textwidth}\centering\small Unencrypted\\snapshots\end{column}
  \begin{column}{0.19\textwidth}\centering\small Monitoring\\logs\end{column}
  \begin{column}{0.19\textwidth}\centering\small CI/CD\\pipeline\end{column}
  \begin{column}{0.19\textwidth}\centering\small Batch\\endpoint\end{column}
  \begin{column}{0.19\textwidth}\centering\small Public\\registry\end{column}
\end{columns}

\end{frame}

% =============================================================================
\section{Implementation}
% =============================================================================

\begin{frame}{Three-Phase Implementation Roadmap}
\note{[3 min] Practical guidance. Phase 1 (weeks 1--4): foundation IT security
adapted for ML (blocks 80\% of commodity attacks). Phase 2 (weeks 5--12): ML-specific
defenses --- DP, watermarking, backdoor scans (5--10\% accuracy cost). Phase 3
(weeks 13--24): advanced --- TEEs, adversarial training, red team (30--100\% extra
training time). Key insight: start with foundation; each phase builds on the previous.
If short: just show the diagram and hit the key numbers.}

% --- Layout: FULL-WIDTH diagram ---
\centering
\safeimg[width=0.92\textwidth,height=4.8cm,keepaspectratio]{implementation-roadmap.pdf}

\vspace{0.1cm}
\small
\mlsysinsight{Phase 1 blocks 80\% of attacks:}{Standard IT security adapted for ML pipelines is the highest-leverage starting point.}

\end{frame}

\begin{frame}{Distributed Systems Multiply the Surface}
\note{[2 min] The shift from 1 machine to N machines increases attack vectors
by approximately N-squared (all-to-all) or N*log(N) (ring). 128 GPUs in 16
machines: compromising one node injects poisoned gradients globally. Edge FL
with 10K clients: 1\% compromise (100 devices) degrades accuracy 10--50\%.
Key insight: attack surface grows superlinearly with scale.
Ask: ``Why is this worse than linear growth?''}

\small
\begin{columns}[T]
  \begin{column}{0.55\textwidth}
    \textbf{Scale multiplies vulnerabilities:}

    \vspace{0.15cm}
    {\footnotesize
    \renewcommand{\arraystretch}{1.2}
    \begin{tabular}{@{}lr@{}}
      \toprule
      \textbf{Topology} & \textbf{Attack Vectors} \\
      \midrule
      All-to-All & $O(n^2)$ \\
      Ring        & $O(n \log n)$ \\
      \bottomrule
    \end{tabular}
    }

    \vspace{0.15cm}
    \begin{itemize}\setlength\itemsep{1pt}
      \item {\footnotesize 128 GPUs / 16 nodes: 1 compromised node poisons the global model}
      \item {\footnotesize FL with 10K clients: 1\% compromise $\to$ 10--50\% accuracy drop}
      \item {\footnotesize Centralized 0.1\% $\equiv$ distributed 0.01\% per node}
    \end{itemize}
  \end{column}
  \begin{column}{0.42\textwidth}
    \vspace{0.1cm}
    \begin{mlsyscard}{errorstroke}
    \textbf{Superlinear growth}\\[0.05cm]
    {\scriptsize Single-node $\to$ distributed ML is a fundamentally expanded attack surface requiring holistic threat modeling.}
    \end{mlsyscard}

    \vspace{0.1cm}
    {\scriptsize\textbf{Required:} Secure inter-node channels, cross-domain identity mgmt, coordinated policies.}
  \end{column}
\end{columns}

\end{frame}

% =============================================================================
\section{Wrap-Up}
% =============================================================================

\begin{frame}{Fallacies}
\note{[2 min] Four common misconceptions. Hit each one with the quantitative
counter-evidence. The key fallacy is ``security through obscurity'' --- black-box
extraction works with 10K--100K queries.}

\scriptsize
\textbf{Fallacy:} \textit{Security through obscurity protects ML models.}\\
Black-box extraction reconstructs 90\% accuracy with 10K--100K queries, within typical API rate limits.

\vspace{0.06cm}
\textbf{Fallacy:} \textit{Differential privacy automatically ensures privacy.}\\
Budgets compound: 10 runs at $\epsilon=1$ each $=$ total $\epsilon=10$. Organizations exceed guarantees by orders of magnitude.

\vspace{0.06cm}
\textbf{Fallacy:} \textit{Federated learning inherently provides privacy.}\\
Membership inference: 70--90\% accuracy on FL models. Gradient inversion reconstructs training data.

\vspace{0.06cm}
\textbf{Fallacy:} \textit{Data poisoning requires compromising large portions of data.}\\
0.1\% (1K in 1M) reduces accuracy 10--50\%. Backdoors need only 0.01\% for 90\%+ attack success.

\end{frame}

\begin{frame}{Pitfalls}
\note{[2 min] Two operational pitfalls. Focus on the first: component-level
security without system-level threat modeling. If short: cover just the first one.}

\footnotesize
\textbf{Pitfall:} \textit{Treating security as an isolated component.}\\
Strong API defenses are bypassed through batch endpoints, monitoring logs, or public registries. A single weak link compromises the entire system. Security is a \textbf{system property}, not a component property.

\vspace{0.15cm}
\textbf{Pitfall:} \textit{Underestimating distributed attack surface expansion.}\\
The shift from 1 to $n$ machines increases attack vectors by $O(n^2)$, not $O(n)$. Inter-node channels, federated aggregation, and edge deployment each create new vulnerability classes requiring dedicated threat modeling.

\end{frame}

% --- RETRIEVAL PRACTICE ---

% --- MUDDIEST POINT ---
\begin{frame}{Muddiest Point}
\note{[2 min] Quick anonymous poll. Students write on a slip of paper or submit
digitally. Collect and scan for patterns. Address the top 2--3 confusions in the
next lecture's opening. This closes the feedback loop.}

\centering
\vspace{1.0cm}
{\Large\bfseries What was the \alert{muddiest point} today?}

\vspace{0.8cm}
{\normalsize Write down the concept you found \textbf{most confusing}.}

\vspace{0.5cm}
{\small\textcolor{midgray}{Anonymous. One sentence. Submit before you leave.}}

\end{frame}

\begin{frame}{What Were the Key Ideas?}
\note{[2 min] Retrieval practice. Students write 90 seconds, no notes.
Do NOT show next slide yet.}

\centering
\vspace{1.5cm}
{\Large\bfseries Close your notes.}

\vspace{0.8cm}
{\large Write down the \textbf{4 most important concepts} from today.}

\vspace{0.8cm}
{\normalsize\textcolor{midgray}{90 seconds --- no peeking.}}

\end{frame}

\begin{frame}{Key Takeaways}
\note{[2 min] Reveal takeaways. Walk through each bullet. Emphasize the quantitative
anchors: 100,000x extraction asymmetry, 0.1\% poisoning leverage, 15\% isolation tax,
epsilon compounding.}

\scriptsize
\begin{itemize}\setlength\itemsep{1pt}
  \item \textbf{Security $\neq$ Privacy}: Security defends against \emph{adversarial manipulation}; Privacy protects against \emph{unauthorized inference}. Both required.
  \item \textbf{Model theft is cheap}: 100,000$\times$ cost asymmetry --- extraction via 10K--100K queries vs.\ \$100M training.
  \item \textbf{Data poisoning has extreme leverage}: 0.1\% corruption $\to$ 10--50\% accuracy drop; backdoors at 0.01\% persist through retraining.
  \item \textbf{Differential privacy is the standard}: $\epsilon$-$\delta$ guarantees beat anonymization, but budgets compound --- track cumulative $\epsilon$.
  \item \textbf{Hardware is the root of trust}: TEEs (5--30\% cost), Secure Boot, HSMs, PUFs anchor software protections.
  \item \textbf{Layered defense}: No single mechanism suffices. Composition creates graceful degradation.
  \item \textbf{Distributed systems multiply the surface}: Attack vectors grow $O(n^2)$, not $O(n)$.
\end{itemize}

\end{frame}

\begin{frame}{References}
\note{[1 min] Point students to key papers.}

\scriptsize
\mlsysref{Farwell+11}{Farwell \& Rohozinski. ``Stuxnet and the Future of Cyber War.'' Survival, 2011.}
\mlsysref{Dwork+06}{Dwork et al. ``Calibrating Noise to Sensitivity in Private Data Analysis.'' TCC, 2006.}
\mlsysref{Abadi+16}{Abadi et al. ``Deep Learning with Differential Privacy.'' CCS, 2016.}
\mlsysref{Goodfellow+14}{Goodfellow et al. ``Explaining and Harnessing Adversarial Examples.'' ICLR, 2015.}
\mlsysref{Biggio+12}{Biggio et al. ``Poisoning Attacks Against SVMs.'' ICML, 2012.}
\mlsysref{Tramer+16}{Tram\`er et al. ``Stealing ML Models via Prediction APIs.'' USENIX Security, 2016.}

\end{frame}

\begin{frame}{Next Lecture: Robust AI}
\note{[1 min] Forward hook. Security addresses intentional threats --- but a model
perfectly protected from attackers can still fail when the real world shifts beneath
it. Chapter 14 shifts focus from malicious threats to operational stress:
distribution drift, hardware faults, compounding failures.}

\footnotesize
\begin{columns}[c]
  \begin{column}{0.30\textwidth}
    \centering
    \begin{mlsyscard}{errorstroke}
    {\large\bfseries Distribution Drift}\\[0.15cm]
    {\footnotesize The world changes; the model doesn't.}
    \end{mlsyscard}
  \end{column}
  \begin{column}{0.30\textwidth}
    \centering
    \begin{mlsyscard}{routingstroke}
    {\large\bfseries Hardware Faults}\\[0.15cm]
    {\footnotesize Silent bit-flips corrupt predictions.}
    \end{mlsyscard}
  \end{column}
  \begin{column}{0.30\textwidth}
    \centering
    \begin{mlsyscard}{computestroke}
    {\large\bfseries Cascading Failures}\\[0.15cm]
    {\footnotesize One failure triggers ten more.}
    \end{mlsyscard}
  \end{column}
\end{columns}

\vspace{0.3cm}
\centering
{\normalsize Security protects against \textbf{malicious} threats.\\
Robustness protects against \textbf{operational} stress.}\\[0.1cm]
{\small\textbf{A defended fleet must also be a resilient fleet.}}

\end{frame}



\appendix

\begin{frame}{Backup: Extended Reference}
\note{Backup slide with additional reference material for this chapter.}

\footnotesize
This slide provides extended reference material for students who want to go deeper.
Refer to the textbook chapter for complete derivations and additional worked examples.

\vspace{0.3cm}
\begin{mlsyscard}{crimson}
\textbf{Key equations and frameworks from this chapter} are collected in the
textbook's summary tables. Use them as a quick reference during problem sets.
\end{mlsyscard}

\end{frame}

\begin{frame}{Backup: Further Reading}
\note{Backup slide. Point students to additional resources beyond the references slide.}

\footnotesize
\textbf{For deeper exploration:}

\vspace{0.15cm}
\begin{itemize}\setlength\itemsep{2pt}
  \item The textbook chapter contains 2--3 additional worked examples
  \item Problem sets apply these concepts to real-world scenarios
  \item The course website links to open-source implementations
  \item Office hours: bring your toughest questions
\end{itemize}

\end{frame}


\end{document}
